%
% 000-session15.tex
% Presentation Session 15: Practice 5 -- BASES2 Web Server (HTTP module)
%
% Compile to .pdf with LaTeX (pdflatex)
% Requires Beamer package (latex-beamer)
%

\documentclass[xcolor=table]{beamer}
\usetheme{Warsaw}
\beamertemplatenavigationsymbolsempty
\setbeamertemplate{headline}{}
\setbeamerfont{title}{size=\large}
\useoutertheme{infolines}
\usepackage[english]{babel}
\usepackage[utf8]{inputenc}
\usepackage{graphics}
\usepackage{amssymb}
\usepackage{multirow}
\usepackage{xcolor}
\usepackage{framed}
\usepackage{hyperref}
\usepackage{tikz}
\usetikzlibrary{arrows.meta, positioning}

\definecolor{shadecolor}{RGB}{180,180,180}
\definecolor{lightgreen}{RGB}{144,238,144}

%%%%%%%%%%%%%%%%%%%%%%%%%%%%%%%%%%%%%%%%%%%%%%%%%%%%%%%%%%%%%%%%
% Python Highlighting
%%%%%%%%%%%%%%%%%%%%%%%%%%%%%%%%%%%%%%%%%%%%%%%%%%%%%%%%%%%%%%%%
\DeclareFixedFont{\ttb}{T1}{txtt}{bx}{n}{10}
\DeclareFixedFont{\ttm}{T1}{txtt}{m}{n}{10}

% Custom colors
\usepackage{color}
\definecolor{deepblue}{rgb}{0,0,0.5}
\definecolor{deepred}{rgb}{0.6,0,0}
\definecolor{deepgreen}{rgb}{0,0.5,0}
\definecolor{grey1}{rgb}{0.5,0.5,0.5}

\usepackage{listings}

% Python style for highlighting
\newcommand\pythonstyle{\lstset{
language=Python,
basicstyle=\ttm\small,
morekeywords={self},
keywordstyle=\ttb\color{deepblue},
emphstyle=\ttb\color{deepred},
stringstyle=\color{deepgreen},
commentstyle=\small\color{grey1}\ttm,
frame=single,
showstringspaces=false
}}

% HTML style for highlighting
\newcommand\htmlstyle{\lstset{
language=HTML,
basicstyle=\ttm\small,
keywordstyle=\ttb\color{deepblue},
stringstyle=\color{deepgreen},
commentstyle=\small\color{grey1}\ttm,
frame=single,
showstringspaces=false
}}

% Python environment
\lstnewenvironment{python}[1][]
{
\pythonstyle
\lstset{#1}
}
{}

% HTML environment
\lstnewenvironment{htmlcode}[1][]
{
\htmlstyle
\lstset{#1}
}
{}

% Python for inline
\newcommand\pythoninline[1]{{\pythonstyle\lstinline!#1!}}

\newcommand{\asignatura}{Session 15: Practice 5 -- BASES2 Web Server}
\newcommand{\grado}{Programming in Network Environments}
\newcommand{\curso}{2025-2026}

\title[\asignatura]{\asignatura}
\subtitle{\grado}
\author{URJC}
\institute{GSyC}
\date{Course \curso}

\AtBeginSection[]
{
\begin{frame}<beamer>
\begin{center}
{\LARGE \insertsection}
\end{center}
\end{frame}
}

\begin{document}

\frame{
\maketitle
}

\begin{frame}
\frametitle{Contents}
\tableofcontents
\end{frame}

%%%%%%%%%%%%%%%%%%%%%%%%%%%%%%%%%%%%%%%%%%%%%%%%%%%%%%%%%%%%%%%%
%%%%%%%%%%%%%%%%%%%%%%%%%%%%%%%%%%%%%%%%%%%%%%%%%%%%%%%%%%%%%%%%
% SECTION 1: Session 14 Review
%%%%%%%%%%%%%%%%%%%%%%%%%%%%%%%%%%%%%%%%%%%%%%%%%%%%%%%%%%%%%%%%
%%%%%%%%%%%%%%%%%%%%%%%%%%%%%%%%%%%%%%%%%%%%%%%%%%%%%%%%%%%%%%%%

\section{Review: Session 14 -- HTTP Python Library}

\begin{frame}[fragile]
\frametitle{What We Learned in Session 14 (1/2)}

\textbf{HTTP Python module}: implement web servers without raw sockets.

\begin{itemize}
\item \textbf{Built-in web server}: \texttt{python -m http.server 8080}
  \begin{itemize}
  \item Serves files from the current directory automatically
  \end{itemize}
\vspace{0.15cm}
\item \textbf{Web Server 1}: Using the \textbf{default handler}
  \begin{itemize}
  \item \pythoninline{http.server.SimpleHTTPRequestHandler}
  \item Works like the built-in server, from Python code
  \end{itemize}
\vspace{0.15cm}
\item \textbf{Web Server 2}: Implementing our \textbf{own handler}
  \begin{itemize}
  \item Create a class derived from \pythoninline{BaseHTTPRequestHandler}
  \item Override the \pythoninline{do_GET()} method
  \end{itemize}
\end{itemize}

\end{frame}

\begin{frame}[fragile]
\frametitle{What We Learned in Session 14 (2/2)}

\textbf{Generating HTTP responses} with the handler:

\vspace{0.1cm}

\begin{python}
class TestHandler(http.server.BaseHTTPRequestHandler):

    def do_GET(self):
        contents = "I am the happy server! :-)"

        self.send_response(200)
        self.send_header('Content-Type', 'text/plain')
        self.send_header('Content-Length',
                         len(contents.encode()))
        self.end_headers()
        self.wfile.write(contents.encode())
\end{python}

\vspace{0.2cm}

\textbf{Key properties}: \pythoninline{self.requestline}, \pythoninline{self.command}, \pythoninline{self.path}

\vspace{0.1cm}

\textbf{Debugging}: Browser developer tools (Network tab)

\end{frame}

%%%%%%%%%%%%%%%%%%%%%%%%%%%%%%%%%%%%%%%%%%%%%%%%%%%%%%%%%%%%%%%%
%%%%%%%%%%%%%%%%%%%%%%%%%%%%%%%%%%%%%%%%%%%%%%%%%%%%%%%%%%%%%%%%
% SECTION 2: Session 15 -- Practice 5
%%%%%%%%%%%%%%%%%%%%%%%%%%%%%%%%%%%%%%%%%%%%%%%%%%%%%%%%%%%%%%%%
%%%%%%%%%%%%%%%%%%%%%%%%%%%%%%%%%%%%%%%%%%%%%%%%%%%%%%%%%%%%%%%%

\section{Session 15: Practice 5 -- BASES2 Web Server}

\begin{frame}
\frametitle{Session 15 Goals}

\textbf{In this session we will}:

\begin{itemize}
\item Reimplement the \textbf{BASES web server} from P04 using the \textbf{HTTP Python module}
\item Use our \textbf{own handler} instead of raw sockets
\item Serve the same DNA base pages: \textbf{A}, \textbf{C}, \textbf{G}, \textbf{T}
\item Handle \textbf{any filename} as an HTML resource
\end{itemize}

\vspace{0.4cm}

\textbf{Time}: 2 hours

\vspace{0.2cm}

\textbf{Folder}: \texttt{P05/} with subfolder \texttt{html/info/}

\vspace{0.2cm}

\textbf{Important}: Finish all S14 exercises before starting!

\end{frame}

\begin{frame}
\frametitle{From P04 to P05: What Changes?}

\textbf{P04} (raw sockets) $\rightarrow$ \textbf{P05} (HTTP module):

\vspace{0.3cm}

\begin{columns}
\begin{column}{0.45\textwidth}
\textbf{P04 -- Raw sockets}:
\begin{itemize}
\item Manual socket creation
\item Parse raw HTTP request
\item Build response string manually (status line, headers, body)
\item \pythoninline{cs.send(response.encode())}
\end{itemize}
\end{column}
\begin{column}{0.45\textwidth}
\textbf{P05 -- HTTP module}:
\begin{itemize}
\item \pythoninline{socketserver.TCPServer}
\item \pythoninline{self.path} gives the resource
\item \pythoninline{self.send_response()}
\item \pythoninline{self.send_header()}
\item \pythoninline{self.wfile.write()}
\end{itemize}
\end{column}
\end{columns}

\vspace{0.4cm}

The \textbf{logic is the same}, but the HTTP module handles the protocol details for us.

\end{frame}

\begin{frame}
\frametitle{BASES2 Server: Resource Table}

The server must handle the \textbf{same resources} as P04:

\vspace{0.2cm}

\begingroup\small
\begin{center}
\begin{tabular}{|l|l|l|}
\hline
\textbf{Resource (path)} & \textbf{HTML File} & \textbf{Color} \\
\hline
\texttt{/info/A.html} & \texttt{html/info/A.html} & lightgreen \\
\hline
\texttt{/info/C.html} & \texttt{html/info/C.html} & yellow \\
\hline
\texttt{/info/G.html} & \texttt{html/info/G.html} & lightblue \\
\hline
\texttt{/info/T.html} & \texttt{html/info/T.html} & pink \\
\hline
\texttt{/} or \texttt{/index.html} & \texttt{html/index.html} & white \\
\hline
any other & \texttt{html/error.html} & red \\
\hline
\end{tabular}
\end{center}
\endgroup

\vspace{0.2cm}

\textbf{Key difference with P04}: Resources now use \texttt{.html} extension (e.g., \texttt{/info/A.html} instead of \texttt{/info/A}).

\end{frame}

\begin{frame}[fragile]
\frametitle{Server Structure with HTTP Module}

\textbf{Server filename}: \texttt{P05/bases2-webserver.py}

\vspace{0.1cm}

\begin{python}
import http.server
import socketserver
from pathlib import Path

PORT = 8080
socketserver.TCPServer.allow_reuse_address = True

class Bases2Handler(
        http.server.BaseHTTPRequestHandler):

    def do_GET(self):
        # 1. Map "/" to "/index.html"
        #    Otherwise: "html" + self.path
        # 2. Open the file (Path.read_text())
        #    If FileNotFoundError: use error.html
        # 3. Send response (status, headers, body)
        pass

Handler = Bases2Handler
with socketserver.TCPServer(("", PORT), Handler) \
        as httpd:
    httpd.serve_forever()
\end{python}

\end{frame}

\begin{frame}
\frametitle{HTML Files: Reusing from P04}

Reuse the HTML files from P04 with \textbf{small modifications}:

\vspace{0.2cm}

\begin{itemize}
\item \texttt{P05/html/index.html} -- Main page (white background)
  \begin{itemize}
  \item Links to \texttt{/info/A.html}, \texttt{/info/C.html}, \texttt{/info/G.html}, \texttt{/info/T.html}
  \end{itemize}
\vspace{0.15cm}
\item \texttt{P05/html/info/A.html} -- Adenine (lightgreen)
\item \texttt{P05/html/info/C.html} -- Cytosine (yellow)
\item \texttt{P05/html/info/G.html} -- Guanine (lightblue)
\item \texttt{P05/html/info/T.html} -- Thymine (pink)
\vspace{0.15cm}
\item \texttt{P05/html/error.html} -- Error page (red background)
\end{itemize}

\vspace{0.3cm}

\textbf{Pay attention}: The info pages are now inside \texttt{html/info/}, so links between pages may need adjusting.

\end{frame}

\begin{frame}
\frametitle{Testing the Server}

\textbf{Steps to test}:

\begin{enumerate}
\item Run: \texttt{python P05/bases2-webserver.py}
\item Open browser: \texttt{http://localhost:8080}
\item Verify the \textbf{main page} appears (index.html)
\item Click on each base link (A, C, G, T)
\item Try a non-existing URL (e.g., \texttt{/doesnotexist})
  \begin{itemize}
  \item The \textbf{error page} should appear
  \end{itemize}
\item Use \textbf{developer tools} (Network tab) to inspect requests
\end{enumerate}

\vspace{0.3cm}

\textbf{Remember}: If you see cached content, press \textbf{F5} or \textbf{Ctrl+Shift+R} to force a reload.

\end{frame}

\begin{frame}
\frametitle{Deliverables}

\textbf{P05 folder should contain}:

\vspace{0.2cm}

\begingroup\small
\begin{itemize}
\item \texttt{bases2-webserver.py} -- The web server using the HTTP module
\item \texttt{html/index.html} -- Main page
\item \texttt{html/error.html} -- Error page
\item \texttt{html/info/A.html} -- Adenine page
\item \texttt{html/info/C.html} -- Cytosine page
\item \texttt{html/info/G.html} -- Guanine page
\item \texttt{html/info/T.html} -- Thymine page
\end{itemize}
\endgroup

\vspace{0.3cm}

\textbf{Remember}: Push everything to your remote Github repository!

\end{frame}

%%%%%%%%%%%%%%%%%%%%%%%%%%%%%%%%%%%%%%%%%%%%%%%%%%%%%%%%%%%%%%%%
%%%%%%%%%%%%%%%%%%%%%%%%%%%%%%%%%%%%%%%%%%%%%%%%%%%%%%%%%%%%%%%%
% SECTION 3: End
%%%%%%%%%%%%%%%%%%%%%%%%%%%%%%%%%%%%%%%%%%%%%%%%%%%%%%%%%%%%%%%%
%%%%%%%%%%%%%%%%%%%%%%%%%%%%%%%%%%%%%%%%%%%%%%%%%%%%%%%%%%%%%%%%

\begin{frame}
\begin{center}
\Huge{Questions?}

\vspace{1cm}

\Large{Time to code!}
\end{center}
\end{frame}

\end{document}
